\section{Conclusion}
\label{sec:conclusion}

For a decade, tabular CFR was the workload that GPU accelerators could not win. GPU-CFR
ends that. It is the first CFR solver on a GPU that beats efficient CPU
implementations, and by a wide margin: \turnSteadyMs\,ms per CFR$^{+}$ iteration on
the 83k-Infoset HUNL turn subgame, \kimGraphRange$\times$ faster than the previous
fastest GPU solver \cite{kim2026parallelizing} on identical A100 hardware, and
\liteGraphLargestRange$\times$ faster than LiteEFG \cite{liu2024liteefg} on the four largest games, with a gap
that grows with the tree. The idea behind the result is general. A fixed game is a
program whose control flow is known before the solver starts, so compiling it once into
flat arrays, folding chance into a template, scheduling by depth, and making every edge
branchless turns a pointer-chasing traversal into 64 to 152 tensor operations that a
CUDA graph replays in a single launch. The pipeline applies to any two-player
zero-sum game, any regret-minimization rule, and any tensor framework, and the iterates
agree bitwise with the reference implementation. Because the compiled sweep is the inner loop of subgame re-solving
\cite{burch2014solving,li2026real}, abstraction refinement
\cite{li2024rl,li2025efficient,li2026effective,li2026abstraction}, and neural CFR
\cite{brown2019deep}, every solver built on tabular CFR inherits the
speedup without changing its algorithm, and offline solves that once needed a cluster
become a job for one GPU. The full source code of GPU-CFR, including the compiler,
the game specifications, and the benchmark scripts, is available at
\url{https://github.com/lbn187/GPU-CFR}, so that the compiled representation can
serve as the starting point for future game solving hardware and software.
